\PassOptionsToPackage{colorlinks=true,allcolors=black}{hyperref}
% REQ: The above line, \PassOptionsToPackage{} must be very first line.

\documentclass[11pt]{article}

% Encoding and font setup for cross-platform reproducibility
\usepackage[T1]{fontenc}
\usepackage{lmodern}

% Improve readability for dense theoretical material
\usepackage{setspace}
\onehalfspacing

% Mathematical notation and table formatting for formal semantics
\usepackage{amsmath, amssymb}
\usepackage{array}
\usepackage[table]{xcolor}
\usepackage{colortbl}
\usepackage{booktabs}
\usepackage{graphicx}

% Theorem environments
\usepackage{amsthm}

% ------------------------------------------------------------
% Numbering convention:
% - All theorem-like statements share one counter, numbered by section.
% - Prose cites statements by stable names and labels, not by number.
% - External documents cite by stable tag (e.g., se100.def.Substrate).
%   Numbers move when text is edited; names do not.
% ------------------------------------------------------------
\theoremstyle{plain}
\newtheorem{theorem}{Theorem}[section]
\newtheorem{lemma}[theorem]{Lemma}
\newtheorem{proposition}[theorem]{Proposition}
\newtheorem{constraint}[theorem]{Constraint}
\newtheorem{claim}[theorem]{Claim}
\newtheorem{assumption}[theorem]{Assumption}

\theoremstyle{definition}
\newtheorem{definition}[theorem]{Definition}
\newtheorem{example}[theorem]{Example}

\theoremstyle{remark}
\newtheorem{remark}[theorem]{Remark}
\newtheorem{note}[theorem]{Note}

% Unnumbered versions for front-matter previews/callouts
\newtheorem*{constraint*}{Constraint}
\newtheorem*{definition*}{Definition}
\newtheorem*{example*}{Example}
\newtheorem*{remark*}{Remark}
\newtheorem*{note*}{Note}

% Bibliography management and hyperlink support
\usepackage{url}
\usepackage{natbib}
\bibpunct{(}{)}{;}{a}{}{ }

% Core hyperlink and bookmark infrastructure
\usepackage{hyperref}
\usepackage{bookmark}

% Support multiple author affiliations
\usepackage{authblk}

% Improve typographic quality and line breaking
\usepackage{microtype}

% Categorical and commutative diagrams
\usepackage{tikz}
\usepackage{tikz-cd}
\usetikzlibrary{arrows.meta,calc,fit,positioning,shapes.geometric}

% Callout boxes for examples and emphasis
\usepackage[most]{tcolorbox}

% Compact list formatting for dense conceptual content
\usepackage{enumitem}
\setlist[itemize]{itemsep=0.2em, topsep=0.2em, parsep=0em, partopsep=0em}
\setlist[enumerate]{itemsep=0.2em, topsep=0.2em, parsep=0em, partopsep=0em}

% Paragraph spacing prioritizes readability over traditional indentation
\setlength{\parskip}{0.75em}
\setlength{\parindent}{0em}

% Keywords macro for structured abstract metadata
\providecommand{\keywords}[1]{\vspace{0.5em}\noindent\textbf{Keywords: }\small #1}

% ------------------------------------------------------------
% Stable named-reference helpers.
% Use labels such as se100.def.Substrate and cite by stable name:
% \DefRef{se100.def.Substrate}{Substrate}
% ------------------------------------------------------------
\newcommand{\DefRef}[2]{\hyperref[#1]{Definition (#2)}}
\newcommand{\AssumpRef}[2]{\hyperref[#1]{Assumption (#2)}}
\newcommand{\ConstRef}[2]{\hyperref[#1]{Constraint (#2)}}
\newcommand{\ExRef}[2]{\hyperref[#1]{Example (#2)}}
\newcommand{\RemRef}[2]{\hyperref[#1]{Remark (#2)}}
\newcommand{\SecRef}[1]{Section~\ref{#1}}

% ------------------------------------------------------------
% Formal notation
% ------------------------------------------------------------
\newcommand{\Substrate}{\mathcal{S}}
\newcommand{\SubstrateRef}{\mathcal{S}_{\mathrm{ref}}}
\newcommand{\Frameworks}{\mathbb{F}}
\newcommand{\Framework}{F}
\newcommand{\FrameworkOne}{F_1}
\newcommand{\FrameworkTwo}{F_2}
\newcommand{\entails}{\vdash}
\newcommand{\notentails}{\nvdash}
\newcommand{\Asserts}{\mathsf{Asserts}}
\newcommand{\pred}[1]{\mathit{#1}}

% Reusable callout box
\newcommand{\FigureCallout}[2]{%
  \begin{tcolorbox}[
    colback=gray!5,
    colframe=black!40,
    title={#1},
    fonttitle=\bfseries,
    arc=3pt,
    boxrule=0.5pt,
    width=\linewidth,
    enhanced,
    breakable
  ]
  #2
  \end{tcolorbox}
}

% =============================================================================
% PREFERRED CITATION SPECIFICATION
%
% @techreport{case2026neutral,
%   author      = {Case, Denise M.},
%   title       = {Neutral Substrates: A Design Constraint for Shared Records Under Persistent Interpretive Disagreement},
%   institution = {Preprint},
%   year        = {2026},
%   url         = {}
% }
% =============================================================================

% ============================================================
\title{
Neutral Substrates:\\
A Design Constraint for Shared Records\\
Under Persistent Interpretive Disagreement
}
% ============================================================

\author{Denise M. Case}
\affil{\small Department of Computer Science and Information Systems\\
Northwest Missouri State University, Maryville, MO, USA\\}
\date{}

\begin{document}
\maketitle
\vspace{-2em}

\begin{abstract}
  Shared accountability records are often used
  by parties who do not, and may never,
  agree about causation, responsibility, or normative interpretation.
  In such settings, a shared record cannot remain neutral
  by omitting contested information,
  because accountability requires preserving what parties
  claimed, believed, inferred, or relied on.
  Nor can it remain neutral
  by asserting one contested interpretation as the shared factual base.

  This paper defines a neutral substrate as a shared representational layer
  that provides stable reference
  while making no substrate-layer causal or normative commitments.
  The central design constraint is that,
  when causal and normative propositions are
  contestable across admissible frameworks and
  the substrate's referential commitments
  are common ground among those frameworks,
  a substrate is robustly neutral by design
  if and only if its foundational layer is restricted to
  framework-invariant referential structure.
  Causal and normative content may still be represented,
  but only as attributed, provenance-bearing assertions
  made by some framework, source, agent, institution, or document.

  The representational machinery used here is standard:
  reification, attribution, and provenance.
  The contribution is the constraint itself:
  a precise, checkable condition on the foundational layer of a shared record,
  stated together with the assumptions it depends on and the boundary condition
  under which it fails.
  A neutral substrate says enough to preserve accountability,
  but it does not turn one party's interpretation into substrate fact.
  Neutrality is unavailable when the referential regime itself is
  contested among the frameworks in play.
\end{abstract}

\keywords{
  Formal ontology;
  neutral substrates;
  accountable records;
  referential regimes;
  interpretive non-commitment;
  extension stability;
  persistent interpretive disagreement;
  reification;
  provenance;
  pre-causal;
  pre-normative
}

\begin{note*}[Citation Convention]
  Definitions, assumptions, constraints, and examples in this paper carry stable
  names (for example, \texttt{se100.def.Substrate}).
  This paper cites statements by stable name rather than by number,
  so citations remain valid even if numbering changes.
\end{note*}

% =============================================================================
% Preview of Main Constraint
% =============================================================================

\bigskip
\noindent\textbf{Preview of Main Constraint.}
\emph{
  Let $\Substrate$ be a substrate
  intended to remain usable by every admissible interpretive framework,
  including admissible frameworks not enumerable when
  $\Substrate$ is designed.
  Suppose causal and normative propositions are
  contestable across admissible frameworks, and
  suppose the referential commitments of $\Substrate$
  are common ground among those frameworks.
  Then $\Substrate$ satisfies interpretive non-commitment and extension stability
  if and only if its foundational layer
  is restricted to framework-invariant referential structure,
  making in particular no substrate-layer commitment
  to any causal or normative proposition.
  A substrate satisfying both properties is neutral.
}

\medskip
\noindent
\emph{Informally.} A neutral substrate may record causal and normative claims,
but only as attributed, provenance-bearing assertions about the shared
referential layer, never as substrate fact.

\section{Introduction}
\label{sec:introduction}

Accountability records are rarely interpreted by only one party.
A record may be used by institutions, auditors, affected individuals,
regulators, courts, researchers, vendors, advocates, or future maintainers.
These parties may agree on some referential facts:
that a record concerns a particular person, event, decision,
instrument, place, file, policy, rule, or institutional artifact.
At the same time, they may disagree, possibly permanently,
about what caused what, who was responsible,
whether a decision was justified, which rule applied,
or which normative standard should control.

This creates a structural design problem.
A shared record must remain usable by multiple interpretive frameworks,
including frameworks that permanently disagree.
If the shared layer makes a substrate-layer commitment
to a contested causal or normative proposition,
the record becomes structurally aligned with one framework against another.
It may still be useful to the favored framework, but it no
longer functions as a neutral substrate for shared accountability.

The purpose of this paper is to define the substrate-level constraint
needed to avoid that failure.
The central claim is simple:

\begin{quote}
  A neutral substrate may record causal or normative claims,
  but it may not make substrate-layer commitments to them.
  It must represent them as attributed, provenance-bearing assertions
  about the referential layer.
\end{quote}

The distinction between asserting a proposition and
representing that some source asserts it is the mechanism of neutrality,
and it is the difference between
a record that suppresses contested content and
a record that preserves contested content without adopting it.

The mechanism is not new (\SecRef{sec:related-work}).
The contribution is the constraint:
a biconditional connecting interpretive non-commitment and extension stability
to a checkable restriction on the substrate's foundational layer,
together with the assumptions and boundary condition on which the constraint depends.

The paper proceeds as follows.
\SecRef{sec:related-work} credits the foundations and states the contribution.
\SecRef{sec:problem-setting} states the problem.
\SecRef{sec:definitions} gives the core definitions.
\SecRef{sec:neutrality} states and justifies the neutrality constraint.
\SecRef{sec:mechanism} explains the reification mechanism.
\SecRef{sec:example} gives a worked example.
\SecRef{sec:boundary} states the boundary condition under which neutrality is unavailable.
\SecRef{sec:non-goals} identifies non-goals.
\SecRef{sec:conclusion} concludes.

\section{Related Work and Contribution}
\label{sec:related-work}

The representational mechanism this paper relies on is not novel.
Prior work developed the machinery for
recording that a source asserts a proposition,
rather than asserting the proposition itself.
This is standard in knowledge representation, logical AI, and argumentation theory.

In knowledge representation on the web,
RDF reification \citep{hayes2004rdf}
and named graphs \citep{carroll2005named}
provide standard machinery often used for
statements about statements, and the
W3C provenance ontology
\citep{lebo2013prov} standardizes attribution of content to
agents, activities,
and sources.
The form $\Asserts(\Framework,\varphi)$ used throughout this paper
is deliberately their form; nothing representational is added.

In logical AI, \citet{mccarthy1993} treats truth as relative to a context via
$\mathit{ist}(c,p)$,
and framework-relative derivability
($\Substrate \cup \Framework \entails p$)
is a restrained cousin of that idea.
The difference is one of purpose:
context logic supports reasoning across contexts,
whereas the substrate defined here deliberately abstains from
cross-framework conclusions at its foundational layer.
Assumption-based truth maintenance \citep{dekleer1986atms}
likewise maintains conclusions relative to
assumption sets and keeps mutually inconsistent worlds simultaneously
available;
extension stability, defined below,
is a record-level analogue of that discipline.
The difference again is one of purpose:
an ATMS is a reasoning architecture,
while the present constraint is a design condition on records.

In argumentation theory, \citet{dung1995} showed how to compute acceptability
among conflicting arguments.
The aim here is complementary and prior:
argumentation frameworks adjudicate among conflicting claims,
whereas a neutral substrate refuses adjudication at the base layer,
preserving the conflicting claims with attribution.
A neutral substrate is the kind of shared input over
which a Dung-style evaluation could later be run by a framework, not by the
substrate.

Two older distinctions underwrite the whole construction.
Speech act theory distinguishes asserting a proposition
from reporting that someone asserted it
\citep{searle1969}.
And the analysis of ontological commitment\citep{guarino1998}
makes precise the idea that a representation commits its
users to a view of the domain;
the substrate-layer commitment defined below is
a record-level application of that idea.

\textbf{The contribution.}
None of these literatures states a checkable
condition on the foundational layer of a shared accountability record.
Each supplies an ingredient:
reification and provenance supply the mechanism;
context logic and truth maintenance supply the relativized-derivability
reading;
argumentation theory supplies the downstream consumer;
speech act theory and ontological commitment supply the underlying distinctions.
The contribution of this paper is the assembly:
an admissibility gate for frameworks,
a contestability assumption for causal and normative content,
a common-ground assumption for reference, and
a pairwise stability condition,
tied together in a single biconditional constraint on the substrate layer,
stated with the honesty conditions that make it usable,
including the exact boundary at which it fails.
Each ingredient is known.
The constraint, and the discipline of saying precisely
when it does not apply, is the addition.

\section{Problem Setting}
\label{sec:problem-setting}

A shared record can fail neutrality in two opposed ways.

First, it can omit accountability-relevant claims entirely.
This avoids substrate-layer commitment,
but it also suppresses the information needed to understand
who claimed what, under what authority, and on what basis.
A record that suppresses contested claims cannot support accountability
for those claims.

Second, it can make a substrate-layer commitment
to a contested causal or normative proposition.
This preserves the content of the claim,
but strips away the attribution that makes the claim contestable.
Such a record may appear complete, but is structurally non-neutral.

A neutral substrate must avoid both failures.
It must preserve contested claims without adopting them
as substrate-layer commitments.
This paper distinguishes three layers:

\begin{enumerate}
  \item \textbf{Referential structure}: entities, occurrences, institutional
        artifacts, identifiers, timestamps, provenance, and the relations needed to
        fix what is being referred to.
  \item \textbf{Attributed assertions}: claims made by frameworks, sources,
        agents, institutions, records, or documents about the referential layer.
  \item \textbf{Substrate-layer commitments}: propositions asserted by the
        substrate itself, independently of any interpretive framework.
\end{enumerate}

Neutrality requires that contested causal and normative propositions remain in
the second layer rather than being promoted to the third.

\section{Definitions}
\label{sec:definitions}

\begin{definition}[Substrate]
  \label{se100.def.Substrate}
  A \emph{substrate} $\Substrate$ is a shared representational base providing
  stable reference for entities, occurrences, and institutional artifacts
  across a set of interpretive frameworks $\Frameworks$.
  Stable reference consists of individuation, co-reference, and persistence.
\end{definition}

A substrate is not the whole record system.
It is the foundational layer on which
interpretive, causal, normative, evidentiary, explanatory, or
institutional claims may be layered.
Its purpose is to make shared reference
possible without forcing agreement about every interpretation of what is
referenced.
A substrate may identify that there is an event, entity, decision,
document, instrument, rule, policy, observation, or institutional artifact, and
it may preserve identifiers, timestamps, provenance, and persistence relations
among them.
Whether some event caused another,
whether some decision was justified, or
whether some rule ought to apply
are not commitments of the substrate
just because the substrate identifies
the relevant event, decision, rule,
or institutional artifact.

\begin{definition}[Substrate-Layer Commitment]
  \label{se100.def.SubstrateCommitment}
  A substrate $\Substrate$ \emph{commits} to a proposition $p$ when
  $\Substrate \entails p$; that is, $p$ is asserted by the substrate
  independently of any interpretive framework.
  A \emph{causal or normative commitment} is a substrate-layer commitment
  to a causal or normative proposition.
  Representing such a proposition as an attributed,
  provenance-bearing assertion does not
  constitute a substrate-layer commitment
  to the asserted proposition:
  \[
    \Substrate \entails \Asserts(\Framework,\varphi)
    \quad\text{does not imply}\quad
    \Substrate \entails \varphi
  \]
\end{definition}

The substrate commits to the attributed assertion
$\Asserts(\Framework,\varphi)$;
it does not commit to $\varphi$ itself.
This distinction is central.
Neutrality does not require suppressing causal or normative information.
It requires keeping causal and normative information at the level
of attributed claims rather than substrate fact.

\begin{definition}[Referential Regime]
  \label{se100.def.ReferentialRegime}
  A \emph{referential regime} is the triple of
  \begin{enumerate}
    \item individuation conditions,
    \item co-reference conditions, and
    \item persistence conditions
  \end{enumerate}
  by which a substrate fixes and tracks entities, occurrences, and institutional
  artifacts.
\end{definition}

Individuation conditions determine when something counts as one thing rather
than another. Co-reference conditions determine when two references refer to
the same thing. Persistence conditions determine when something remains the
same thing across time, transformation, revision, or institutional change. The
referential regime is what allows a substrate to support shared use: if parties
cannot agree on what is being referred to, they cannot reliably disagree about
causes, norms, explanations, or obligations concerning that thing.

\begin{definition}[Admissible Framework]
  \label{se100.def.AdmissibleFramework}
  A framework $\Framework$ is \emph{admissible} if it satisfies three conditions:
  \begin{enumerate}
    \item \textbf{Internal consistency}: $\Framework \notentails \bot$.
    \item \textbf{Empirical honesty}: where $\Framework$ makes empirical claims,
          it cites evidence and does not contradict scientific consensus in the
          relevant domain.
    \item \textbf{Documented interpretive function}: $\Framework$ is presented as
          an interpretive function with a
          named source, documented scope, and citable basis.
  \end{enumerate}
\end{definition}

Admissibility is not agreement.
Two admissible frameworks may disagree about
causation, responsibility, justification, applicability, or normative
significance.
The point of admissibility is not to decide which framework is
correct.
It is to distinguish frameworks that can participate in a shared
accountability record
from arbitrary, inconsistent, fabricated, or purely
strategic rationalizations.

The admissibility gate is not itself neutral,
and this paper does not claim that.
Condition~2 invokes scientific consensus,
and the identification of consensus is a judgment
that can be contested at the meta-level.
Condition~3, as stated, distinguishes documented frameworks from undocumented ones,
not sincere frameworks from strategic ones:
a well-documented post-hoc rationalization satisfies the provenance requirement.
The gate's actual
function is accordingly weaker and more honest than legitimacy screening.
It excludes frameworks that cannot participate in a shared record at all such as
inconsistent, evidence-free, or unattributable ones, and it forces every
admitted framework to be visible, scoped, and citable,
so that even a rationalization enters the record as an
identifiable, examinable position rather than an anonymous one.
The neutrality defined in this paper is therefore conditional:
it is neutrality among admissible frameworks,
relative to an
admissibility judgment that is itself normatively loaded.
That judgment is,
like any other interpretive act,
representable in a record as an attributed
claim with provenance rather than as substrate fact.

\begin{definition}[Framework-Variant Proposition]
  \label{se100.def.FrameworkVariant}
  A proposition $p$ is \emph{framework-variant} with respect to substrate
  $\Substrate$ if there exist admissible frameworks $\FrameworkOne$ and
  $\FrameworkTwo$ such that
  \[
    \Substrate \cup \FrameworkOne \entails p
    \quad\text{and}\quad
    \Substrate \cup \FrameworkTwo \entails \neg p
  \]
\end{definition}

Framework variance captures disagreement that arises not from different
referential bases,
but from different admissible interpretations layered on top
of the same referential base.
A proposition that is stable under one framework
and rejected under another is not suitable as an unqualified substrate-layer
commitment.

\begin{definition}[Framework-Invariant Proposition]
  \label{se100.def.FrameworkInvariant}
  A proposition $p$ is \emph{framework-invariant} with respect to $\Frameworks$
  if, for every admissible framework $\Framework \in \Frameworks$,
  \[
    \Framework \cup \{p\} \notentails \bot
  \]
  that is, $p$ is compatible with every admissible framework.
\end{definition}

\begin{definition}[Framework-Invariant Commitment Set]
  \label{se100.def.FrameworkInvariantCommitmentSet}
  A set of substrate-layer commitments $C$ is \emph{framework-invariant} with
  respect to $\Frameworks$ if, for every admissible framework
  $\Framework \in \Frameworks$,
  \[
    C \cup \Framework \notentails \bot .
  \]
\end{definition}

Invariance is a positive guarantee, not the mere absence of variance. A
proposition that no admissible framework derives, but that some admissible
framework refutes, is neither framework-variant (so long as nothing asserts it)
nor framework-invariant.
The neutrality constraint below requires invariance:
compatibility with every admissible framework, including admissible frameworks
that join $\Frameworks$ after the substrate is designed.

\begin{assumption}[Contestability]
  \label{se100.def.Contestability}
  Let $C_{cn}$ be the class of causal or normative propositions
  whose interpretation lies within the contested domain.
  No $p \in C_{cn}$ is guaranteed framework-invariant at design time.
\end{assumption}

Contestability does not mean that all causal or normative propositions are
false, arbitrary, or unknowable.
It means that the substrate cannot certify
universal compatibility for such propositions across the admissible frameworks
for which the record must remain usable.
The assumption is domain-relative:
it applies where records are expected to support
accountability under persistent disagreement
about causation, norms, responsibility, legitimacy, or institutional interpretation.

\begin{assumption}[Referential Common Ground]
  \label{se100.def.ReferentialCommonGround}
  Let $\SubstrateRef$ denote the \emph{referential commitments} of $\Substrate$.
  These are the substrate-layer commitments fixed by its referential regime:
  identifiers;
  the typing of entities, occurrences, and institutional artifacts;
  timestamps; provenance; and referential relations among them.
  For every admissible framework $\Framework$,
  \[
    \SubstrateRef \cup \Framework \notentails \bot
  \]
  and this compatibility is preserved when $\Frameworks$ is extended by further
  admissible frameworks.
\end{assumption}

\AssumpRef{se100.def.ReferentialCommonGround}{Referential Common Ground}
states that $\SubstrateRef$ is a framework-invariant commitment set and
formalizes the scope condition on which neutrality
depends:
the frameworks in play contest interpretation, not reference.
It says that the substrate's referential commitments are framework-invariant
as a set, not merely one by one.
It is an assumption, not a fact of nature, and it can fail;
\SecRef{sec:boundary} treats its failure, which is exactly the condition under
which neutrality is unavailable.

\begin{definition}[Interpretive Non-Commitment]
  \label{se100.def.InterpretiveNonCommitment}
  A substrate $\Substrate$ satisfies \emph{interpretive non-commitment} if it
  makes no substrate-layer commitment to any framework-variant proposition.
  Equivalently, if
  $p$ is framework-variant with respect to $\Substrate$, then
  $\Substrate \notentails p$ and $\Substrate \notentails \neg p$.
\end{definition}

Interpretive non-commitment does not prevent the substrate from recording that
a framework, source, agent, institution, or document asserts $p$.
It prevents the substrate itself from asserting $p$ independently of attribution.

\begin{definition}[Extension Stability]
  \label{se100.def.ExtensionStability}
  A substrate $\Substrate$ satisfies \emph{extension stability} if, for every
  admissible framework $\Framework$,
  \[
    \Substrate \cup \Framework \notentails \bot
  \]
\end{definition}

This is pairwise compatibility between the substrate and each admissible
framework, not global consistency across all admissible frameworks.
The distinction matters: admissible frameworks may contradict one another.
A neutral substrate need not make all frameworks mutually consistent.
It must remain separately usable by each admissible framework without itself causing
contradiction.

\begin{remark}[Relation Between the Two Properties]
  \label{se100.remark.PropertyRelation}
  Extension stability entails interpretive non-commitment.
  The two properties are retained under separate names
  because they identify distinct design failures:
  commitment on disputed content and breakdown of layerability.
\end{remark}

\section{The Neutrality Constraint}
\label{sec:neutrality}

\begin{constraint}[Robust Neutrality]
  \label{se100.constraint.Neutrality}
  Let $\Substrate$ be a substrate intended
  to remain usable by every admissible framework,
  including admissible frameworks not enumerable when $\Substrate$ is designed.
  Assume
  \AssumpRef{se100.def.Contestability}{Contestability} and
  \AssumpRef{se100.def.ReferentialCommonGround}{Referential Common Ground}.
  Then $\Substrate$ is robustly neutral by design
  iff its foundational layer is restricted
  to the referential commitments $\SubstrateRef$;
  in particular, $\Substrate$ makes no substrate-layer commitment
  to any causal or normative proposition.
\end{constraint}

Informally: if a shared record must remain usable by admissible frameworks that
may permanently disagree about causation and norms, then the shared substrate
may assert only reference.
Causal and normative content must be represented as
attributed, provenance-bearing claims
about the referents fixed by the substrate, not as substrate-layer commitments.

A substrate satisfying interpretive non-commitment and extension stability is neutral.
A substrate satisfying those properties
by the referential restriction above is robustly neutral by design.

This is a design constraint and the justification is short.
It is stated in full so that its dependence on the
two assumptions is clear.

\emph{Sufficiency.} Suppose the foundational layer of $\Substrate$ is
restricted to $\SubstrateRef$. By
\AssumpRef{se100.def.ReferentialCommonGround}{Referential Common Ground},
$\Substrate \cup \Framework \notentails \bot$ for every admissible $\Framework$,
so extension stability holds; by
\RemRef{se100.remark.PropertyRelation}{Relation Between the Two Properties},
interpretive non-commitment follows.
Because referential common ground is
preserved under admissible extension of $\Frameworks$,
the guarantee does not
depend on enumerating the frameworks in advance.
Note that sufficiency uses the
common-ground assumption essentially:
referential content is not inherently
uncontestable, and where that assumption fails,
so does this direction
(\SecRef{sec:boundary}).

\emph{Design-time necessity.} Suppose $\Substrate$ makes a
substrate-layer commitment to a proposition $p$ outside $\SubstrateRef$.
If $p$ is causal or normative and lies within the contested domain,
\AssumpRef{se100.def.Contestability}{Contestability} says that $p$ is not
guaranteed framework-invariant at design time.
Since $\Substrate$ is intended to remain usable across admissible frameworks
not enumerable when $\Substrate$ is designed,
$\Substrate$ cannot certify extension stability while making that commitment.

If $p$ is neither causal, normative, nor referential,
then $p$ falls outside $\SubstrateRef$
and outside the contested causal-or-normative class $C_{cn}$.
Its compatibility with every admissible framework
therefore remains uncertified at design time.
The one principled exception is treated in
\RemRef{se100.remark.ConsensusProtected}{Consensus-Protected Content}, and it
marks the edge of the constraint's necessity.

\begin{remark}[Consensus-Protected Content]
  \label{se100.remark.ConsensusProtected}
  The empirical-honesty condition of
  \DefRef{se100.def.AdmissibleFramework}{Admissible Framework} carves out one
  class of non-referential content that is framework-invariant by construction:
  propositions whose negation would contradict settled scientific consensus in
  the relevant domain.
  Strictly, a substrate could make substrate-layer commitments to such
  propositions without violating extension stability,
  because no admissible framework may refute them.
  The conservative policy is to attribute even these propositions to
  the measuring instrument, the standards body, or the scientific record
  for two reasons.
  First, consensus shifts over the lifetime of
  accountability records, and a substrate-layer commitment that was protected
  at design time may not remain so.
  Second, the identification of consensus is itself a
  judgment, made by someone, at some time, and on some basis.
  That is exactly the profile of content this paper places in the attributed layer.
  The referential restriction in
  \ConstRef{se100.constraint.Neutrality}{Neutrality} is therefore the stated
  design policy;
  this remark marks the one place where its necessity is
  epistemic and prudential rather than logical.
\end{remark}

Robust neutrality at design time is achieved by separating layers.

\section{Mechanism: Reification of Interpretive Claims}
\label{sec:mechanism}

The mechanism of neutrality is reification.
A non-neutral substrate asserts a
contested causal or normative proposition:
\[
  \Substrate \entails \varphi
\]
A neutral substrate instead represents
that some framework, source, agent,
institution, or document
asserts the proposition:
\[
  \Substrate \entails \Asserts(\Framework,\varphi)
\]
The substrate commits to the attributed assertion
$\Asserts(\Framework,\varphi)$,
not to the asserted proposition $\varphi$
(\DefRef{se100.def.SubstrateCommitment}{Substrate-Layer Commitment}).

This allows the record to preserve accountability-relevant content without
making the shared layer partisan among admissible frameworks.
The substrate may
record the decision that was made; the person or entity affected;
the time of the decision;
the instrument or system involved;
the policy or rule cited;
the source that produced an explanation;
the text of that explanation;
the causal or normative proposition the source asserts;
and the provenance of the assertion.
The substrate must not convert the asserted explanation into a
substrate-layer commitment to the asserted proposition.

\section{Worked Example}
\label{sec:example}

\begin{example}[Reification Fragment]
  \label{se100.example.ReificationFragment}
  Consider a decision record concerning an automated eligibility denial. The
  record involves a subject $u$, a decision $d$, a model or instrument $m$, a
  timestamp $t$, an institutional policy $r$, and an explanation produced by
  source $\Framework$.
\end{example}

The example presupposes referential common ground:
the frameworks in play agree
which subject, which decision, which instrument, and which policy the record
concerns.
What they contest is interpretation.

\textbf{Non-neutral form.} The substrate makes substrate-layer commitments
to causal or normative propositions:
\[
  \Substrate \entails \pred{caused\_by}(\pred{score}(u,m), \pred{denial}(d))
  \qquad
  \Substrate \entails \pred{justified\_by}(\pred{denial}(d), \pred{policy}(r))
\]
If admissible frameworks disagree about whether the score caused the denial,
whether the policy applied, or whether the denial was justified,
these commitments make the substrate non-neutral.

\textbf{Neutral form.} The substrate asserts the referential structure:
\begin{align*}
  \Substrate & \entails \pred{Subject}(u)
             & \Substrate                    & \entails \pred{occurred\_at}(d,t)     \\
  \Substrate & \entails \pred{Decision}(d)
             & \Substrate                    & \entails \pred{involved}(d,u)         \\
  \Substrate & \entails \pred{Instrument}(m)
             & \Substrate                    & \entails \pred{used\_instrument}(d,m) \\
  \Substrate & \entails \pred{Policy}(r)
             & \Substrate                    & \entails \pred{cites}(d,r)
\end{align*}
and records the contested content as attributed claims:
\[
  \Substrate \entails \Asserts(\Framework,\;
  \pred{caused\_by}(\pred{score}(u,m), \pred{denial}(d)))
\]
\[
  \Substrate \entails \Asserts(\Framework,\;
  \pred{justified\_by}(\pred{denial}(d), \pred{policy}(r)))
\]
The substrate preserves that source $\Framework$
made these causal and normative assertions.
It commits to the attributed assertions, not to the asserted propositions.

\textbf{Predicate-level reading.} The neutral record contains the event,
actors, timestamps, identifiers, institutional artifacts, and
provenance-bearing claims.
It contains no unattributed top-level content whose
meaning is equivalent to the non-neutral assertions above.
This holds regardless of field naming:
a field named \texttt{reason}, \texttt{cause},
\texttt{basis}, \texttt{risk}, \texttt{fault}, or \texttt{violation}
may look like descriptive metadata
while functioning as a substrate-layer
commitment to a causal or normative proposition.
Causal or normative vocabulary may appear;
the test is whether the substrate makes a substrate-layer commitment
to a causal or normative proposition.

The neutral form supports
multiple downstream interpretations without
requiring the substrate to choose among them.
An institutional framework may treat the explanation as valid.
An audit framework may question whether the cited policy actually applied.
A claimant-side framework may contest the causal explanation.
A regulatory framework may evaluate whether the explanation satisfies disclosure requirements.
The substrate remains usable by all of them
because it preserves the shared referential base and the attributed assertions
without making a substrate-layer commitment
to the contested interpretation.

\section{Boundary Condition: When Neutrality Is Unavailable}
\label{sec:boundary}

Neutrality depends on
\AssumpRef{se100.def.ReferentialCommonGround}{Referential Common Ground}.
If admissible frameworks disagree not merely about causal or normative
interpretations, but about the substrate's individuation, co-reference, or
persistence conditions,
then some admissible framework is incompatible with $\SubstrateRef$ itself.
The assumption fails, and neutrality is unavailable at that layer.

For example, one framework may treat two records as referring to the same
institutional event while another admissible framework treats them as distinct
events. Or one framework may treat an institutional identity as persistent
across a merger while another treats the merger as identity-breaking.
In such cases, the disagreement concerns the referential regime itself, and the
substrate cannot be neutral among frameworks that do not share enough
referential structure to identify what is being interpreted.

This does not make accountability impossible.
It means the contested
referential regime must itself be represented as an attributed claim, or moved
to a higher-level structure capable of tracking the disagreement.
What the original substrate cannot do is declare one referential regime neutral while
admissible frameworks contest the individuation, co-reference, or persistence
conditions on which that declaration depends.

This boundary condition prevents the constraint from claiming neutrality
where reference itself is contested.
Without it, neutrality could become a rhetorical device for laundering a contested
referential scheme into the shared base.
A neutral substrate must not hide normative or causal commitments inside reference.

\section{Scope and Non-Goals}
\label{sec:non-goals}

\textbf{This paper does not propose a data format.} A neutral substrate may be
implemented using relational databases, graph structures, event logs, RDF-like
representations, document stores, schemas, or application-specific formats.
The neutrality constraint applies to the commitments made by the substrate layer,
not to any serialization or storage technology.

\textbf{This paper does not adjudicate which frameworks are legitimate in any
  particular case.} Admissibility identifies minimum conditions for participation
in the shared record: consistency, empirical honesty, and provenance.
It does not decide which admissible framework should prevail in a dispute.

\textbf{This paper does not claim a neutral admissibility gate.} As stated in
\SecRef{sec:definitions}, the gate embeds judgments, including what counts as
consensus and what counts as a documented interpretive function.
The neutrality defined here is conditional on those judgments.
The gate's role is to make participation visible and citable, not to certify sincerity.

\textbf{This paper does not solve accountability generally.} Neutrality is a
structural condition for shared records under persistent disagreement.
It does not determine liability, responsibility, institutional legitimacy, causal
truth, moral blame, legal compliance, or policy adequacy.
It makes such disputes representable without forcing the shared substrate to decide them
prematurely.

\textbf{This paper does not prohibit causal or normative content.} Under
the contestability assumption, it prohibits substrate-layer commitments
to causal or normative propositions in the foundational layer.
Causal and normative claims remain fully representable as attributed,
provenance-bearing assertions.

\section{Discussion}
\label{sec:discussion}

The neutrality constraint reframes neutrality as a structural property of
shared records rather than as a claim about objectivity, consensus, or absence
of values.
A record is not neutral because it contains no contested material;
such a record would often be inadequate for accountability. Nor is a record
neutral because it contains every party's explanation in a single
undifferentiated field; such a record may preserve text while still failing to
distinguish substrate fact from attributed interpretation.

Neutrality requires a specific representational discipline:
the substrate may
assert reference, but it must attribute interpretation.
This discipline matters
most in systems that combine records from multiple institutions, jurisdictions,
tools, or governance contexts, where causal and normative terms often appear
deceptively stable as ordinary metadata.
Neutral design does not prohibit such terms;
it asks what the substrate is committing to
(\DefRef{se100.def.SubstrateCommitment}{Substrate-Layer Commitment}).
If the substrate asserts that a denial was caused by a score, that a policy justified
a decision, or that a party violated a norm, the substrate has taken an
interpretive position.
If it records that a specified source made that claim
under documented provenance, the substrate has preserved the claim without
adopting it.

The constraint has practical value only where referential common ground holds.
The paradigm case is adversarial record examination,
including litigation, forensic discovery, audit, and regulatory review.
In these settings, parties agree which record, which decision,
which subject, and which instrument are at issue while disputing,
often permanently, what caused the outcome and whether it was justified.
This is the common condition of records under dispute.
The hard accountability cases in which reference itself is contested,
including disputed event identity and contested institutional persistence, are real.
They are exactly the boundary cases of \SecRef{sec:boundary}.
This paper claims applicability where common ground holds and unavailability where it does not.
It does not assume reference is always shared.

\section{Conclusion}
\label{sec:conclusion}

A shared accountability record must often serve parties who disagree
permanently about causation, norms, responsibility, legitimacy, or
institutional interpretation.
Such a record cannot remain neutral if its
foundational layer asserts one contested interpretation as fact.

The neutral substrate constraint is therefore straightforward:

\begin{quote}
  The shared layer may assert framework-invariant referential structure.
  Causal and normative content must be represented as
  attributed, provenance-bearing claims, not as substrate fact.
\end{quote}

This does not make the record empty.
It makes the record usable for accountability.
The substrate can preserve who said what, about which entity, event, decision,
instrument, policy, or artifact, under what provenance, and within what scope.
The substrate must not convert an attributed interpretation into a
substrate-layer commitment.

Neutrality is unavailable when the referential regime itself is contested.
But where enough reference can be shared, a neutral substrate provides the
structural base for disagreement without erasure, accountability without
premature adjudication, and interpretation without hidden commitment.

\section*{Statements and Declarations}

\subsection*{Author Contributions}
The author is solely responsible for the
conception, analysis, and writing of this work.

\subsection*{Declaration of Conflicting Interest}
The author declares no potential conflicts of interest with respect to the
research, authorship, and/or publication of this work.

\begin{thebibliography}{99}

  \bibitem[Carroll et~al.(2005)]{carroll2005named}
  Carroll, J.~J., Bizer, C., Hayes, P., and Stickler, P. (2005).
  Named graphs, provenance and trust.
  In \emph{Proceedings of the 14th International Conference on World Wide Web
    (WWW 2005)}, pages 613--622. ACM.

  \bibitem[de~Kleer(1986)]{dekleer1986atms}
  de~Kleer, J. (1986).
  An assumption-based TMS.
  \emph{Artificial Intelligence}, 28(2):127--162.

  \bibitem[Dung(1995)]{dung1995}
  Dung, P.~M. (1995).
  On the acceptability of arguments and its fundamental role in nonmonotonic
  reasoning, logic programming and $n$-person games.
  \emph{Artificial Intelligence}, 77(2):321--357.

  \bibitem[Guarino(1998)]{guarino1998}
  Guarino, N. (1998).
  Formal ontology and information systems.
  In \emph{Formal Ontology in Information Systems (FOIS~'98)}, pages 3--15.
  IOS Press.

  \bibitem[Hayes(2004)]{hayes2004rdf}
  Hayes, P. (2004).
  \emph{RDF Semantics}.
  W3C Recommendation, World Wide Web Consortium.

  \bibitem[Lebo et~al.(2013)]{lebo2013prov}
  Lebo, T., Sahoo, S., and McGuinness, D.~L., editors (2013).
  \emph{PROV-O: The PROV Ontology}.
  W3C Recommendation, World Wide Web Consortium.

  \bibitem[McCarthy(1993)]{mccarthy1993}
  McCarthy, J. (1993).
  Notes on formalizing context.
  In \emph{Proceedings of the 13th International Joint Conference on Artificial
    Intelligence (IJCAI-93)}, pages 555--560.

  \bibitem[Searle(1969)]{searle1969}
  Searle, J.~R. (1969).
  \emph{Speech Acts: An Essay in the Philosophy of Language}.
  Cambridge University Press.

\end{thebibliography}

\end{document}
